\documentclass[main.tex]{subfiles}
\begin{document}

\color{black}

Những ngày tháng tại căn mancave của \hls\ vẫn tiếp tục trôi qua, và dường như càng làm việc lâu với các anh tài, Kaiho càng nhận ra rằng công việc tổng quản của mình chẳng bao giờ chỉ gói gọn trong những nhiệm vụ thông thường. Từ những buổi luyện tập nghiêm túc cho đến các trò đùa chẳng biết bắt nguồn từ đâu, mỗi ngày đều có thể mở ra một tình huống hoàn toàn mới. Và trong mười tuần tiếp theo, cô sẽ còn phải chứng kiến nhiều điều kỳ lạ hơn nữa!

Trước hết, Miyabi, Temma và Aruran được giao nhiệm vụ luyện tập cho một buổi meet and greet, dưới sự giám sát của Kaiho và Ryuuhou. Tuy nhiên, việc chuẩn bị để gặp gỡ người hâm mộ không đơn giản chỉ là học cách chào hỏi hay trò chuyện, nhất là khi những người tham gia đều có cách thể hiện bản thân rất riêng. Trong một câu chuyện khác, Miyabi lại bất ngờ tự nhốt mình trong phòng, khiến Shien, Temma, Kaiho và Suzuno phải nghĩ đủ mọi cách để đưa cậu ra ngoài. Sau nhiều nỗ lực tưởng chừng vô ích, lời giải cho vấn đề hóa ra lại nằm ở một thứ vô cùng đơn giản: chiếc bụng đang đói của Miyabi.

Những trò đùa trong căn mancave cũng không thiếu những màn tranh luận chẳng đi đến đâu. Oga và Rikka có một trận vật tay, nhưng điều đáng chú ý hơn cả lại là màn đổi tên của hai người. Chứng kiến chuyện ấy, Izuru, Astel và Roberu cũng bắt đầu bàn luận về tên gọi, khiến một cuộc thi đấu thể lực vô tình biến thành buổi tranh luận ngôn từ. Còn Astel thì tiếp tục mang đến những câu chuyện về người ngoài hành tinh, với những lời kể khiến người nghe không biết nên tin tưởng, nghi ngờ hay đơn giản là để cậu nói cho hết.

Bên cạnh đó, kỹ năng giao tiếp đầy tính áp đặt của Oga cũng được đem ra thử nghiệm. Từ cách nói chuyện có phần độc đoán ấy, cậu còn tìm cách huấn luyện những người xung quanh trở thành bản sao của mình. Thế nhưng, đằng sau màn đào tạo kỳ quặc đó lại là một sự thật khác: tất cả chỉ là nội dung của một tựa game! Cũng trong những khoảnh khắc sinh hoạt tưởng chừng bình thường, một bài thơ lục bát được viết, với sự góp mặt của cả những người đã ngủ gật lẫn những người vẫn còn thức. Một sự kết hợp giữa thi ca và tình trạng thiếu ngủ, tạo nên những vần thơ mà có lẽ chẳng nhà thơ nào muốn nhận là tác phẩm của mình.

Chưa dừng lại ở đó, Oga, Astel, Temma và Rikka còn hợp sức kể chuyện cười nhằm rèn luyện khả năng nhịn phản bác (tsukkomi) của Roberu. Đáng tiếc, những câu chuyện được đưa ra lại chẳng đủ sức thuyết phục để tạo nên một buổi luyện tập đúng nghĩa. Người phải chịu đựng tất cả không chỉ có Roberu mà còn có cả Kaiho, trong khi Kuon cũng chẳng thể giấu nổi vẻ ngán ngẩm khi chứng kiến màn thử thách này. Có lẽ, đôi khi việc giữ im lặng trước những câu chuyện vô lý còn khó hơn cả việc tìm ra một câu trả lời hợp lý cho chúng.

Một trong những sự kiện đáng nhớ khác bắt nguồn từ ý tưởng Astel muốn xăm Aruran. Từ một đề nghị tưởng chừng chỉ mang tính đùa cợt, mọi chuyện nhanh chóng bị đẩy sang một hướng hoàn toàn khác, dẫn đến cả một màn đám ma bi hài. Và khi trò đùa vượt quá giới hạn chịu đựng, Aruran cuối cùng cũng không thể tiếp tục giữ bình tĩnh, khiến những người còn lại phải đối diện với cơn giận của ông chú trong nhóm. Quả nhiên, không phải trò đùa nào cũng có thể kết thúc bằng tiếng cười.

Trong khi đó, một sự cố điện khác lại xảy ra tại căn mancave, lần này thủ phạm chính là Miyabi khi cậu làm nổ tung cầu chì để tiết kiệm điện, khiến toàn bộ nơi làm việc chìm vào bóng tối. Thay vì chờ đợi mọi thứ được khắc phục, các anh tài lại tranh thủ biến khoảng thời gian mất điện thành một dịp để bày đủ trò nghịch ngợm. Có vẻ như sau những gì từng xảy ra trong lần mất điện trước, bóng tối đã không còn là trở ngại đối với họ nữa, mà trái lại, nó còn trở thành điều kiện lý tưởng để những ý tưởng khó hiểu nhất được tự do phát triển.

Và để khép lại chuỗi câu chuyện lần này, Roberu đã đem đến một bộ phim lấy bối cảnh miền Viễn Tây, xoay quanh mối ân oán đầy éo le của\dots\ một đĩa chiếc bánh kếp. Trong thế giới ấy, Okayu của \hll\ xuất hiện với vai trò chủ quán saloon, còn Kaiho, Suzuno và Ryuuhou cũng tham gia vào câu chuyện. Từ bầu không khí cao bồi cho đến những mối thù chẳng giống ai, bộ phim đã biến một miền Viễn Tây vốn quen thuộc thành sân khấu cho một cuộc đối đầu mà có lẽ chẳng ai có thể đoán trước được kết cục.

Mười tuần tiếp theo lại khép lại với vô số câu chuyện đáng nhớ, từ những buổi luyện tập gặp gỡ người hâm mộ, các màn tranh luận về tên gọi, những bài thơ ngẫu hứng cho đến các trò đùa vượt khỏi giới hạn thông thường. Kaiho tiếp tục chứng kiến những khía cạnh mới của từng thành viên \hls, đồng thời cũng có thêm nhiều dịp phối hợp với Ryuuhou, Suzuno và cả Kuon trong những tình huống bất đắc dĩ. Còn các anh tài \hls\ vẫn giữ nguyên phong độ của mình: không thiếu tài năng, nhưng cũng chẳng bao giờ thiếu những ý tưởng khiến người khác phải đau đầu.

Ba quyển sách đã ghi lại một chặng đường ngày càng dài của \textit{stars*collection}, nơi Kaiho không còn là một tổng quản mới toanh như những ngày đầu tiên, mà đã dần trở thành một phần quen thuộc trong cuộc sống của các anh tài. Dẫu vậy, những câu chuyện tại căn mancave vẫn chưa có dấu hiệu kết thúc. Khi các thành viên tiếp tục làm việc, vui đùa và đồng hành cùng nhau, những tình huống mới chắc chắn sẽ còn xuất hiện, mang theo cả tiếng cười lẫn những phen khiến vị tổng quản phải tự hỏi mình đang làm việc ở đâu.

Hãy cùng đón chờ những câu chuyện tiếp theo trong **Quyển số Bốn: Khóa Lầy lội cao cấp**, nơi Kaiho và các anh tài \hls\ sẽ tiếp tục bước vào những cuộc phiêu lưu mới, với những màn phối hợp ăn ý\dots\ hoặc hoàn toàn trái ngược với điều đó!

\end{document}